% SPDX-License-Identifier: GPL-3.0-or-later
% Copyright (C) 2026 Alessandro Greco
\chapter{Observing Traffic: Wireshark and Logs}\label{cap:20-traffico}
\aggiornatoal{23 September 2026}{tor 0.4.9.13; tshark (Debian trixie); laboratory network}

\begin{obiettivi}
By the end of this chapter, you will be able to:
\begin{itemize}[nosep]
  \item capture Tor traffic with \cmd{tcpdump} and analyze it using \cmd{tshark} or Wireshark;
  \item distinguish what an external observer can see from what remains hidden;
  \item decrypt the TLS on \emph{your own} channel using session keys and inspect its cells (Laboratory~12);
  \item follow the protocol step by step using tor's logs.
\end{itemize}
\end{obiettivi}

In the preceding chapters, we examined the protocol \enquote{from the inside} by writing code to speak it. Here, we observe it \enquote{from the outside}, as someone watching the network would. This exercise serves two opposite purposes: understanding what an adversary can infer from traffic and obtaining an independent way to verify our own implementation.

\begin{attenzione}[Only Your Own Traffic, and Only in the Laboratory]
Capturing other people's traffic without authorization is unlawful in Italy, as in most countries (\cref{cap:25-legale}), and conflicts with the ethics of anonymity research (\cref{sec:rp-perche}). All captures in this chapter involve the book's private network or connections that we open ourselves.
\end{attenzione}

% ---------------------------------------------------------------------------
\section{What an Observer Can See}\label{sec:traffico-osservatore}

Someone observing the connection between a client and its guard, such as an Internet provider, can see:
\begin{itemize}
  \item IP addresses and ports: the client connects to an address that, unless it belongs to a bridge, appears in the public consensus and is therefore recognizable as a Tor relay;
  \item the TLS handshake: certificates travel encrypted under TLS~1.3, but some characteristics of the client's initial message remain visible;
  \item the sizes and timing of TLS records.
\end{itemize}
The observer cannot see cell contents, circuits, or destinations. Yet sizes and timing are precisely what \emph{website fingerprinting} and correlation attacks rely on (\cref{cap:26-attacchi}). Encrypted content does not mean that traffic reveals nothing.

In \cref{sec:lab-censura}, we compared ordinary Tor traffic, recognizable as TLS, against obfs4 traffic with no recognizable structure. Here, we go one step further and look inside TLS.

% ---------------------------------------------------------------------------
\section{The Tools}\label{sec:traffico-strumenti}

\begin{description}
  \item[\cmd{tcpdump}] Captures packets from a network interface and saves them to a \texttt{pcap} file. For example, inside the laboratory container, capture traffic to \texttt{test004r}'s ORPort:
\begin{lstlisting}[language=bash,numbers=none]
$ docker exec -u 0 torbook-network tcpdump -i lo \
    -w /lab/net/tls.pcap "tcp port 5104"
\end{lstlisting}
  \item[Wireshark and \cmd{tshark}] Analyze packet captures, the former through a graphical interface and the latter from the command line. They recognize TLS, but there is no official Tor-protocol dissector. Even after decryption, Tor cells appear as raw bytes. Our \file{cells.py} module will serve as the dissector.
\end{description}

\begin{figure}[htbp]
\centering
% SPDX-License-Identifier: GPL-3.0-or-later
% Copyright (C) 2026 Alessandro Greco
%% Accessibility alternative text, used only in EPUB.
\altfig{Simplified view of Tor over TCP and TLS. External observers can see IP addresses, ports, traffic volumes, timing and TLS record lengths but not Tor cells without TLS session keys. Multiple cells may be carried in a TLS record, and a cell may be split across records, so record boundaries need not align with cell boundaries. With link protocol v4 or higher, fixed-length cells are 514 bytes, consisting of a 5-byte header and a 509-byte body; variable-length cells also exist. Layered circuit encryption applies to the bodies of RELAY cells, not to all cell types such as PADDING and VPADDING. TLS keys reveal cell framing but do not automatically decrypt RELAY cell bodies.}
% Encapsulation overview (TLS record and Tor cell boundaries need not match).
\begin{tikzpicture}[font=\sffamily\scriptsize,
  liv/.style={draw,rounded corners=2pt,minimum height=0.85cm,align=center},
  cel/.style={draw,fill=torviola!12,minimum height=0.7cm,minimum width=1.5cm,font=\sffamily\tiny}]
  \node[liv,fill=torgrigio!10,minimum width=10.6cm] (tcp) at (0,2.6) {TCP byte stream: variable segments};
  \node[liv,fill=torverde!12,minimum width=9.4cm] (tls) at (0,1.35) {TLS: record / cell borders differ};
  \foreach \i/\x in {1/-3.6,2/-1.8,3/0,4/1.8,5/3.6} \node[cel] (c\i) at (\x,0) {514\,B};
  \node[anchor=west,font=\sffamily\tiny,text=torgrigio] at (4.6,0) {fixed v4+};
  \node[liv,fill=avviso!10,minimum width=4.6cm] (rel) at (-1.7,-1.4) {RELAY: 509\,B body\\layer-encrypted};
  \node[liv,fill=torgrigio!8,minimum width=4.6cm] (pad) at (3.3,-1.4) {PADDING / VPADDING\\not RELAY};
  \draw[torgrigio] (c1.south west) -- (rel.north west); \draw[torgrigio] (c3.south east) -- (rel.north east);
  \node[anchor=west,align=left,font=\sffamily\tiny,text=torgrigio] at (5.8,2.6) {IP, ports, volume\\and timing visible};
  \node[anchor=west,align=left,font=\sffamily\tiny,text=torgrigio] at (5.8,1.35) {TLS record lengths\\visible on the link};
  \node[anchor=west,align=left,font=\sffamily\tiny,text=torgrigio] at (5.8,0) {cells visible after\\TLS decryption};
  \node[anchor=west,align=left,font=\sffamily\tiny,text=torgrigio] at (5.8,-1.4) {RELAY body needs\\circuit keys};
  \node[font=\sffamily\tiny,text=torgrigio,align=center] at (0,-2.55) {Schematic only: TLS record and cell boundaries need not align.};
\end{tikzpicture}

\caption{Conceptual Tor channel: TCP segments carry TLS records whose boundaries need not match cell boundaries. IP addresses, ports, record lengths and timing are observable; TLS decryption exposes cells. Under link protocol v4+, a fixed-length cell is 514 bytes (5-byte header, 509-byte body), though variable-length cells also exist. Only \texttt{RELAY} cell bodies are encrypted in circuit layers; padding cells are not RELAY messages.}
\label{fig:incapsulamento}
\end{figure}

\subsection{Decrypting Your Own TLS}

TLS uses session keys that external observers do not possess. A party \emph{participating} in the connection, however, can export those keys. Browsers and applications using Python's \texttt{ssl} library can export these keys if they enable key logging. Our \file{channel.py} explicitly sets \texttt{SSLContext.keylog\_filename} when \texttt{SSLKEYLOGFILE} is defined. Wireshark and \cmd{tshark} use that file to decrypt a capture. Laboratory~1's \file{channel.py} follows this convention: set the variable before running it. The file contains lines such as:
\begin{lstlisting}[numbers=none]
# TLS secrets log file, generated by OpenSSL / Python
SERVER_HANDSHAKE_TRAFFIC_SECRET 4878224d3f9fb9aedb96d... 0da89dff6feb...
\end{lstlisting}
These record the client random value and the secrets from each phase of TLS~1.3 for each session.

\begin{approfondimento}[Why This Is Not an Attack on Tor]
Decrypting TLS on your own channel reveals nothing that the client did not already know. Even with TLS session keys, \cella{RELAY} cells remain encrypted in layers with keys for individual hops (\cref{cap:10-cifratura-relay}). Reading those would require circuit keys held only by the client and the relays involved. TLS is the onion's outer shell, not its only layer.
\end{approfondimento}

% ---------------------------------------------------------------------------
\section{Laboratory: Cells Inside TLS}\label{sec:lab-traffico}

\begin{laboratorio}[12]{Capturing, decrypting, and reading a channel}
\textbf{Objective.} Capture a channel opened from our client to a relay, decrypt it with session keys, and reconstruct the sequence of cells using \file{labs/en/code/torlab/traffic.py}.\\
\textbf{Prerequisites.} Laboratories~1--3; the laboratory network running; \cmd{tcpdump} and \cmd{tshark} installed in the container.\\
\textbf{Estimated time.} 1--2 hours.
\end{laboratorio}

For Laboratory~12, start the private network and open two terminals at the repository root. Run the capture command below on the \emph{host}: it starts \cmd{tcpdump} as root \emph{inside the laboratory container}, since the ordinary Makefile shell runs under an unprivileged user. Keep the capture running while the client executes in the second terminal, then stop it with Ctrl+C. The file is written to the shared \file{labs/en/lab/net/} volume.

\paragraph{Step 1: capture.} Start capturing inside the container. Then, from the host, open a channel to \texttt{test004r} and build a one-hop circuit while exporting the session keys:
\begin{lstlisting}[language=bash,numbers=none]
$ PYTHONPATH=labs/en/code SSLKEYLOGFILE=labs/en/lab/net/tls-keys.log python3 -c "
from torlab.laboratory import relay
from torlab.circuit import build
build([relay('004r')], cc=False)"
\end{lstlisting}

Run the following Python command from the \emph{repository root} on Linux with host networking. \texttt{PYTHONPATH} makes \file{torlab} importable from \file{labs/en/code/}; \file{labs/en/code/torlab/channel.py} explicitly reads \texttt{SSLKEYLOGFILE} and enables TLS secret logging. The key file contains sensitive session secrets: do not publish it, and delete it after the experiment.

\paragraph{Step 2: identify our connection.} Other laboratory relays also connect to the same port, but their session keys are different. Ask \cmd{tshark} to follow each TCP stream after attempted decryption, retaining the one that decrypts successfully:
\begin{lstlisting}[language=bash,numbers=none]
$ docker exec torbook-network tshark -r /lab/net/tls.pcap \
    -o tls.keylog_file:/lab/net/tls-keys.log \
    -q -z follow,tls,raw,2
\end{lstlisting}
The TCP stream identifier \texttt{2} is only an example from an earlier capture. Discover candidate stream IDs using \cmd{tshark} with \texttt{-T fields -e tcp.stream}, then follow a decryptable stream. The output below is historical; certificate lengths and circuit IDs can change between runs.
The output contains the decrypted data from both directions in hexadecimal; one direction is indented with a tab.

\paragraph{Step 3: inspect cells.} The \texttt{cells\_from\_follow} function reconstructs both streams and passes them to Laboratory~2's \texttt{decode}. Each direction has its own version negotiation: its first cell is \cella{VERSIONS} with a two-byte CircID. Four-byte CircIDs are used only after negotiation. The result:
\begin{lstlisting}[numbers=none]
client -> relay  CircID=0x0000 VERSIONS body=4 bytes
relay  -> client CircID=0x0000 VERSIONS body=6 bytes
relay  -> client CircID=0x00000000 CERTS body=1443 bytes
relay  -> client CircID=0x00000000 AUTH_CHALLENGE body=36 bytes
relay  -> client CircID=0x00000000 NETINFO body=509 bytes
client -> relay  CircID=0x00000000 NETINFO body=509 bytes
client -> relay  CircID=0xc244f1f6 CREATE2 body=509 bytes
relay  -> client CircID=0xc244f1f6 CREATED2 body=509 bytes
\end{lstlisting}
Observed from outside, this is exactly the sequence in \cref{fig:sessione} and \cref{cap:07-canali}. The client's \cella{VERSIONS} cell has a four-byte body because \file{channel.py} offers versions 4 and 5 only. Circuit ID \texttt{0xc244f1f6} has its most significant bit set to 1 because the client opened the channel (\cref{sec:circid}).

\begin{attenzione}[Two Errors in Our First Version]
Our initial analysis made two mistakes, revealed by comparing the output against the expected sequence. It used one negotiated protocol version for both directions, misinterpreting the relay's \cella{VERSIONS} cell. It also assigned blocks to the wrong direction because we guessed which direction \cmd{tshark} had indented. The correct rule is simple: the first block is always the client's \cella{VERSIONS} cell.
\end{attenzione}

% ---------------------------------------------------------------------------
\section{tor Logs}\label{sec:traffico-log}

A packet capture shows \emph{what} travels over the network. tor's logs explain \emph{why}. At \texttt{info} level, tor records circuit construction stages, guard choices, and descriptor downloads and publication. At \texttt{debug}, its reports extend to individual cells. We have already used logs to verify which HSDir the client queried (\cref{sec:lab-onion}) and when an onion service rotated its descriptors.

A few practical suggestions:
\begin{itemize}
  \item \texttt{debug} produces enormous amounts of text. Enable it only for a particular domain, for example \texttt{Log [circ]debug file circ.log}, to limit detailed messages to circuit events.
  \item On the real network, \opt{SafeLogging} masks addresses. A client's logs should never contain the destinations it visits.
  \item Public-relay logs may contain other people's personal data, with the legal consequences discussed in \cref{cap:25-legale}; retain only what is strictly necessary.
\end{itemize}

% ---------------------------------------------------------------------------
\section*{Summary}
\begin{itemize}
  \item A connection observer sees addresses, the TLS handshake, packet sizes, and timing, but not Tor cells or destinations.
  \item \texttt{SSLKEYLOGFILE} permits decryption of TLS on your own channel and inspection of its cells; \cella{RELAY} cells nevertheless remain encrypted in layers.
  \item Analyzing a packet capture provides an independent way to verify an implementation.
  \item tor's \texttt{info} and \texttt{debug} logs explain the program's choices; during real-world use, logging should be minimized and protected.
\end{itemize}

% ---------------------------------------------------------------------------
\section*{Exercises}

\begin{esercizio}[1]\label{es:20-dimensioni}
In a capture of a Tor download, TLS records are not all multiples of 514 bytes long. Why, if cells have a fixed size?
\end{esercizio}

\begin{esercizio}[2]\label{es:20-relay}
Extend Laboratory~12 by decrypting the \cella{RELAY} cells on your own circuit using keys from \file{circuit.py}. Which relay messages can you see while constructing a three-hop circuit?
\end{esercizio}

\begin{esercizio}[2]\label{es:20-padding}
Capture one minute of a laboratory client's channel without generating application traffic. Count TLS records in each direction and connect your observations to channel padding in \cref{cap:07-canali}.
\end{esercizio}

\begin{esercizio}[3]\label{es:20-dissector}
Write a small Wireshark dissector in Lua that recognizes Tor cells within decrypted data and displays their CircID and command.
\end{esercizio}

% ---------------------------------------------------------------------------
\section*{Further Reading}
\begin{itemize}
  \item Wireshark documentation on decrypting TLS using session key log files.
  \item The tor manual's \opt{Log} entry, explaining logging domains and verbosity levels.
\end{itemize}
